\section{Interpretation, limitations and disposition}
\label{a2:limits}

The original manuscript, the F nuisance companion and the present
follow-ups address related but distinct questions. The original observed
recipe selected \(r=0.4\) at its first origin and \(r=0.7\) later, with a
different training magnitude-rate estimator. The new comparisons fix
\(r=0.7\) in their reference fits and use a shared training-only
conditional-within-bin estimator. Their matched contrasts are valid
within those separately frozen protocols; they are not corrected entries
in the original paper's predictive table \citep{ProjectOriginal}.

The F companion still has 1,528 unknown method outcomes out of 8,596
planned. Its historical counts, continuation counts and accounting bounds
retain their original scope. Neither observed follow-up supplies the
missing null-query results, establishes empirical null size or upgrades a
power comparison. No F fit, bank, calibration or resampling plan was
reopened for this editorial addition \citep{ProjectFirstAddendum}.

The class-history construction is deliberately restricted. The consulted
MDFHP preprint permits class-specific temporal and magnitude-productivity
parameters and bounded class mark distributions, with upper magnitude
10 \citep[Section~2]{DavisBaeumerWangPreprint2024}. Our common-kernel,
low-dimensional \(P(h)\) and source-fertility
contrast, together with an unbounded excess-exponential working law, are
different modelling choices. The conditional summed-intensity identity
depends on those choices. It does not refute the broader class-history
framework, reproduce the authors' Japan or Middle America fits, or
transfer identification and stationarity theorems to a nonlinear
inhibitory family.

All forecast windows had been inspected in earlier research. Training
prefixes overlap and the sensitivity settings reuse the same 612
targets. The results are retrospective exploratory procedure comparisons,
not prospective confirmation or independent replications. Three windows
do not support treating clustered events as independent draws for
standard errors or significance claims. Native magnitude rounding,
mixed magnitude conventions, time-varying detection, missing smaller or
deeper parents, pre-1960 activity and national spatial pooling remain
observation/model assumptions. They were not resolved by these
extensions.

The engineering controls and complete numerical audits make the
implementation inspectable, but do not establish global optima, precise
identification, confidence coverage or power. Positive control examples
show that the fitted pipeline can return gains on particular generating
histories. Equal Perron branching radius did not equate immigrant-weighted
offspring production in the class-active control; different offspring
composition and event counts limit any information-matched comparison.
Sparse profiles and near-training-equivalent candidate spreads retain
their descriptive meanings.

The negative forecast results do not establish that physical inhibition
is absent. Nor do lower fitted slopes establish that \(r=0.7\) is
physically wrong, or show that the earlier negative fertility contrasts
were caused by that fixed value. The cross-run pattern motivates a
possible compensation interpretation, while joint estimation of the
fertility/association extension with free \(r\) remains unperformed.
The persistence of \(E0\)'s forecast advantage over \(I0\) is specific
to the tested recipes and windows.

Both investigations are therefore retained at their respective negative
checkpoints. The fixed-\(r\) \(E0\) candidate remains the predictive
reference for these follow-up comparisons. Broken-slope, saturation,
spatial-locality and independent-data extensions were outside their
executed scope. A future study would require a separately justified
question, information set and protocol, rather than further tuning of
these already examined forecast windows.
